\honest{The Martial Art of Rationality}{Eliezer Yudkowsky}{2006}

I often say that rationality is the martial art of the mind. You do not need big muscles to learn a martial art; if you have a hand, you can learn to make a fist. Likewise, if you have a brain, you might learn to use it properly. I admit that athletic people are more likely to take up martial arts, but that ``may be a matter of enjoyment as much as anything else.'' I grant that fast learners learn faster, ``but the art of rationality isn't about that.'' \nb{That is all the post says about the obvious objection, that good reasoning may be mostly talent.}

Our hands obey us better than our minds because control of muscles is old in evolution and reasoning about our own reasoning is new. ``It is not by bigger muscles that the human species rose to prominence upon Earth.'' \nb{That is true, and it is about the species, not about you.}

Then my question. There is a martial arts school in every city, so why is there no school of rationality? \nb{Schools have taught logic for centuries, and statistics for about a century. In 1980 California's state universities began requiring every student to take a course in critical thinking. The post does not mention any of this.}

One reason, perhaps, is that skill in rationality is hard to verify. A student of Tae Kwon Do breaks a board, and everyone sees it. If the fist is wrong, the teacher makes a correct one for the student to copy. \nb{The post does not say what a teacher of reasoning could hold up for a student to copy.}

Martial arts techniques have been refined for generations, and techniques of rationality are harder to pass on. \nb{Logic and mathematics have been refined for generations too.}

In the last few decades we have learned a great deal about human reasoning: the research on heuristics and biases, Bayesian statistics, evolutionary psychology, social psychology. Earlier I named one finding in passing, ``an insensitivity to scope.'' Here I name none, and call them ``new focusing lenses through which to view the landscape of our own minds,'' through which we may see ``the muscles of our brains, the fingers of thought as they move.'' ``Humanity may finally be ready'' to found the martial art of the mind.

I learned what I know about rationality from working on artificial intelligence. Building a mind is harder than training one in most ways, and easier in some. In a human you cannot rewire the machinery, and some of the machinery is built to rationalize.

I admit that building a personal art out of abstract science may be ``awkward,'' like inventing a martial art from physics and anatomy. My answer is that we have an inner eye: introspection works, though it sees blurrily, with systematic distortions. I give no evidence for this. \nb{Nisbett and Wilson had reported in 1977 that people often cannot report the causes of their own judgments, and give confident wrong answers instead.}

I end with three sentences of ``must.'' We must move our own mental limbs, we must connect theory to practice, and we must see what the science means for our daily inner life.
